\chapter{M003 --- Termination Enumeration}

\section{Stage 2 --- Mathematical Analysis}

\subsection{Objective}

Establish the mathematical properties of the transformation from a
Primitive Surface to its finite Termination Set independently of the
implementation.

\subsection{Definition}

A \emph{Termination Set} is the finite collection of physically distinct
surface terminations associated with a primitive surface after
identifying gauge-equivalent representatives.

\subsection{Proposition M003.1}

Every admissible primitive surface admits a finite termination set.

\subsection{Proof Strategy}

The proof proceeds by showing:

\begin{enumerate}
\item admissible truncations are generated from the periodic crystal;
\item gauge-equivalent truncations are quotiented out;
\item physically distinct truncations remain;
\item the resulting set is finite.
\end{enumerate}

\subsection{Open Mathematical Questions}

\begin{itemize}
\item Treatment of polar terminations.
\item Symmetry-equivalent terminations.
\item Stoichiometric versus non-stoichiometric terminations.
\end{itemize}

\subsection{Completion Criterion}

Stage 2 completes once the above propositions are proved without
reference to the implementation.
